\section{Finite quantitative source-chart orders for the retained third seed}
\label{qext:setup}
Fix one retained band, with
\[
0<Q\le1,\quad\varepsilon=Q^h,\quad A=\frac12+h,\quad D=\frac12-h,
\qquad r=\sqrt Q\,R,\quad z=Q^D Z .
\]
Here \(A_0>0\), \(\lambda _0\), and \(\mu=\lambda _0\) retain their original
values.  The source clock \(\nu=\sqrt{A_0}\) is distinct from the physical
viscosity \(\nu_{\rm NS}>0\).  The actual third seed and target are
\[
S_3=\frac{A_0}{\mu}\widehat\lambda _3^{-k_3-6},\quad
\Theta_3(t_3^{\rm in})=-S_3,\quad
-\Theta_3(t_3^a)=S_3e^{L_3}
=\frac{A_0}{\mu}\widehat\lambda _3^{-7/8},\quad
L_3=(k_3+\tfrac{41}{8})\log\widehat\lambda _3 .
\]
The last equality follows by adding the exact exponents
\(-k_3-6+k_3+41/8=-7/8\).  All constants below are finite maxima along
the actual \((\Theta_3,\Omega_3,\zeta_3)\) coefficient path and its actual
parent matrices; no endpoint or sign substitution is made.

\subsection{Physical curl and axial derivative}
For each retained harmonic let
\[
t_j=\chi_jB_jL_jy_3,\quad
c_j=\frac{i\,n_{\Phi,j}\times t_j}{k_{\beta_j}m_j|n_{\Phi,j}|^2},\quad
b_j=t_j+\operatorname{curl}_{*,\mathrm{coeff}}c_j,
\]
where the coefficient curl includes every cutoff, auxiliary chain,
cylindrical \(R^{-1}\), and evolving-path derivative.  Put
\[
v_j=\Re(Q^{-A}b_je^{ik_{\beta_j}m_j\Phi_j}),\quad
d_j=\partial_{Z_j}b_j+ik_{\beta_j}m_j(\partial_{Z_j}\Phi_j)b_j .
\]
For \(e=1,2\), put \(\alpha_e=(e+1)/4-A\). With \(\|a\|_e^2=\int R^e\langle|a|^2\rangle\,dR\),
\[
\|v_j\|_e\le Q^{(e+1)/4-A}\|b_j\|_e,\qquad
\|\partial_zv_j\|_e\le Q^{(e+1)/4-A-D}\|d_j\|_e .
\]
The powers are exact chart powers; the real-part fields obey these finite
upper bounds (a single pure nonzero harmonic has the sharper factor
\(2^{-1/2}\)).  Thus
\[
\begin{array}{c|cc}
e&\text{factor multiplying }\|b_j\|_e&\text{factor multiplying }\|d_j\|_e\\ \hline
2&Q^{1/4-h}&Q^{-1/4}\\
1&Q^{-h}=\varepsilon^{-1}&Q^{-1/2}
\end{array}
\]
as the factors multiplying the actual coefficient norms in the preceding bounds, before summing the finite retained harmonics.

\subsection{Flux products and moments}
Write separate old-parent orders
\[
\|w_{*,\tau}\|_e\le Q^{\omega_{e,\tau}}W_{e,\tau},\qquad
\|\partial_zw_{*,\tau}\|_e\le Q^{\dot\omega_{e,\tau}}\dot W_{e,\tau},
\quad \tau\in\{z,\theta\}.
\]
Put \(B_{j,e,\tau}=\|b_{j,\tau}\|_e\) and
\(D_{j,e,\tau}=\|d_{j,\tau}\|_e\).  Cauchy--Schwarz in the original
\(r^e dr\,d\theta\,dY\) measure gives the separate old--new terms
\[
\begin{aligned}
|J_{2,\mathrm{old}\times v}|&\le Q^{\alpha_2}\bigl(Q^{\omega_{2,z}}W_{2,z}B_{j,2,\theta}
 +Q^{\omega_{2,\theta}}W_{2,\theta}B_{j,2,z}\bigr),\\
|M_{2,\mathrm{old}\times v}|&\le Q^{\dot\omega_{2,z}+\alpha_2}\dot W_{2,z}B_{j,2,\theta}
 +Q^{\omega_{2,z}+\alpha_2-D}W_{2,z}D_{j,2,\theta}\\
&\quad +Q^{\dot\omega_{2,\theta}+\alpha_2}\dot W_{2,\theta}B_{j,2,z}
 +Q^{\omega_{2,\theta}+\alpha_2-D}W_{2,\theta}D_{j,2,z},\\
|J_{1,\mathrm{old}\times v}|&\le 2Q^{\alpha_1+\omega_{1,z}}W_{1,z}B_{j,1,z},\\
|M_{1,\mathrm{old}\times v}|&\le 2Q^{\dot\omega_{1,z}+\alpha_1}\dot W_{1,z}B_{j,1,z}
 +2Q^{\omega_{1,z}+\alpha_1-D}W_{1,z}D_{j,1,z}.
\end{aligned}
\]
The factors (1,2) are finite majorants from
\(|\operatorname{Re}z|\le|z|\); for a single pure harmonic they may be
replaced by the exact Gram factors.  The new quadratic terms are
\[
|J_{2,\mathrm{quad}}|\le Q^{2\alpha_2}B_{j,2,z}B_{j,2,\theta},\quad
|M_{2,\mathrm{quad}}|\le Q^{2\alpha_2-D}(D_{j,2,z}B_{j,2,\theta}+B_{j,2,z}D_{j,2,\theta}),
\]
\[
|J_{1,\mathrm{quad}}|\le Q^{2\alpha_1}B_{j,1,z}^2,\qquad
|M_{1,\mathrm{quad}}|\le 2Q^{2\alpha_1-D}B_{j,1,z}D_{j,1,z}.
\]
Substituting \(A=\tfrac12+h\), \(D=\tfrac12-h\), the quadratic powers are
\[
J_2:Q^{1/2-2h},\quad M_2:Q^{-h}=\varepsilon^{-1},\qquad
J_1:Q^{-2h}=\varepsilon^{-2},\quad M_1:Q^{-1/2-h}.
\]
These are finite norm majorants, not asymptotic source theorems; old-parent
orders remain the explicit \(\omega\) and \(\dot\omega\) symbols above.
The old field and old derivative orders stay independent in all four cross terms. No equal-order cross-term table is inferred from the quadratic powers.

The evaluated-phase correction remains the exact map
\[
M_e^{\rm ev}-M_e=\partial_z\int r^e
 \mathcal S^\circ_{z,\tau_e}(r,z,t,Y(r,t))\,dr,\qquad
 \langle\mathcal S^\circ\rangle_Y=0 .
\]
Define the finite normalized derivative majorant
\[
C_e^\circ(Z,T)=\int_0^\infty R^e\sup_Y
 \left|\partial_Z\bigl[Q^{2A}(\mathcal S^\circ_{z,\tau_e})^Q\bigr](R,Z,T,Y)\right|\,dR .
\]
The direct physical change of variables gives the exact finite envelope
\[
|M_e^{\rm ev}-M_e|\le Q^{(e+1)/2-2A-D}C_e^\circ(Z,T)
 =Q^{e/2-2A+h}C_e^\circ(Z,T).
\]
The zero auxiliary mean therefore does not cancel the evaluated flux.  This
uses the already differentiated normalized chart stress and introduces no
unsupported product of two stress norms.

\subsection{Retained viscosity}
The viscous term is the physical \(\nu_{\rm NS}\Delta v_j\).  In this chart
its radial, axial and angular costs carry \(Q^{-1}\), \(Q^{-2D}=Q^{-1+2h}\),
and \(Q^{-1}R^{-2}\), respectively, including the cylindrical
\(2R^{-2}\partial_\theta v_r\) coupling.  Thus no source \(\varepsilon\)
or clock \(\nu=\sqrt{A_0}\) is substituted for \(\nu_{\rm NS}\).  These
finite identities propagate the actual signed third seed and coefficient
path through \(v,\partial_zv,J_2,J_1,M_2,M_1\), while making no assertion
about a uniform \(Q\downarrow0\) endpoint.
\label{qext:conclusion}
